% ================= APPENDIX: ELECTRICAL SYSTEM DETAIL =================
\section{Electrical System Detail}
\label{app:electrical}

This appendix holds the electrical detail kept out of
Section~\ref{sec:electrical}: the schematic sheets indexed in
Table~\ref{tab:schematic_sheets}, the connector and pin assignments, the harness
drawings, the board layout and the bring-up checklists. The architectural
decisions --- two rails behind one E-stop, linear bus topology, perfboard rather
than fabricated board --- remain in Section~\ref{sec:electrical}. The part
specifications are in Appendix~\ref{app:bom} and are not repeated.

\subsection{Schematic Sheets}
\label{app:schematics}

\subsubsection{Sheet 1 --- Power Entry and Protection}

% Figure placeholder: power entry schematic. XT90 pack connector, XT90-S
% anti-spark inrush path, combined E-stop and inline fuse, split to motor bus
% and logic rail, 470-1000 uF / >= 50 V bulk cap on the pack bus.
% Verify against: fuse rating vs. worst-case bus current
% (Table~\ref{tab:power_budget}); inrush into the driver bulk capacitance --
% the anti-spark path is what prevents welding a bare connector.

\subsubsection{Sheet 2 --- Logic Rail}

% Figure placeholder: LM2596 buck converter 22 V -> 5 V, 1N5819 back-feed
% protection diode, 100 uF/16 V rail bulk capacitance, 100 nF per-device
% decoupling.
% Verify against: 2-3 A converter rating (derates without heatsink) vs. the
% 5 V worst-case total; capacitor voltage ratings -- the 16 V parts belong on
% the logic rail ONLY and must never be fitted to the ~25 V pack bus.

\subsubsection{Sheet 3 --- CAN-FD Bus}

% Figure placeholder: linear bus chain across the three moteus-n1 drivers and
% the Teensy, with split termination (2 x 60.4 ohm + 4.7 nF) at the two
% physical ends.
% Verify against: the topology rules of Table~\ref{tab:signal_rules} --
% linear chain, stubs < 10 cm, termination at the two ends only.

\subsubsection{Sheet 4 --- Encoder Interfaces}

% Figure placeholder: MA600 SPI interface, one per axis, driver to encoder.
% Verify against: SPI run < 20 cm and air gap <= 0.5 mm with the sensor
% centred on the magnet thickness; these fix the driver positions
% (Sec.~\ref{sec:bus_topology}), not the reverse.

\subsubsection{Sheet 5 --- Flight Computer and IMU}

% Figure placeholder: MCU, IMU interface, orientation relative to the body
% frame of Sec.~\ref{sec:cad_method}, and the mounting location relative to the
% pivot required by Sec.~\ref{sec:control2d}.

\subsubsection{Sheet 6 --- Telemetry Link}

% Figure placeholder: telemetry module, serial interface to the MCU, and its
% supply decoupling sized for the transmit burst.

\subsubsection{Sheet 7 --- Brake Servo Drive}

% Figure placeholder: servo supply, decoupling and PWM routing.
% Verify against: stall-current headroom on the logic rail; PWM routed clear of
% the encoder SPI.

\subsection{Connector and Pin Assignments}
\label{app:pinout}

Table~\ref{tab:connectors} is the connector schedule. It exists so that a
harness can be built and re-built from this document alone, and so that a
mis-mate is a documented impossibility rather than a discovered one: no two
connectors carrying different voltages share a housing type on this machine.

\begin{table}[h]
\centering
\caption{Connector schedule. Keying and housing types are chosen so that two
connectors carrying different voltages cannot be mated.}
\label{tab:connectors}
\small
\setlength{\tabcolsep}{4pt}
\begin{tabular}{c L{3.0cm} L{3.0cm} L{2.4cm} L{3.2cm}}
\toprule
Ref & From & To & Housing / keying & Signals \\
\midrule
J1 & Battery pack        & Protection / E-stop   & XT90 & Pack $+$, pack $-$ \\
J2 & Protection          & Distribution board    & XT90 & Motor bus \\
J3 & Distribution board  & Driver, axis 1        & XT30 & Motor bus \\
J4 & Distribution board  & Driver, axis 2        & XT30 & Motor bus \\
J5 & Distribution board  & Driver, axis 3        & XT30 & Motor bus \\
J6 & Board               & CAN chain, end A      & JST-PH3 & CAN-H, CAN-L, GND \\
J7 & CAN chain, end B    & Termination network   & JST-PH3 & CAN-H, CAN-L, GND \\
J8 & Driver, axis $i$    & Encoder, axis $i$     & \textcolor{red}{TBD} & SPI, \qty{3.3}{\volt}, GND \\
J9 & Board               & IMU                   & Qwiic & I\textsuperscript{2}C, supply, GND \\
J10 & Board              & Telemetry module      & \textcolor{red}{TBD} & UART, supply, GND \\
J11 & Board              & Brake servo, axis $i$ & \qty{2.54}{\milli\metre} header & PWM, \qty{5}{\volt}, GND \\
\bottomrule
\end{tabular}
\end{table}

Note the connector mismatch at J2--J5: the pack terminates in XT90 while the
drivers accept XT30, so a distribution harness must be fabricated rather than
bought. It is the item H3 of Table~\ref{tab:wire_schedule} and it is on the
critical path for first power-on.

\subsection{Harness Drawings and Wire Schedule}
\label{app:harness}

Table~\ref{tab:wire_schedule} is the wire schedule. Conductor gauge on the motor
bus is set by the spin-up current rather than the balancing current, since
spin-up is the sustained high-current phase; the signal runs are set by length
limits rather than by current.

\begin{table}[h]
\centering
\caption{Wire schedule. Motor-bus gauge is sized on spin-up current; signal runs
are length-limited (Table~\ref{tab:signal_rules}).}
\label{tab:wire_schedule}
\small
\setlength{\tabcolsep}{4pt}
\begin{tabular}{c L{3.4cm} c c L{3.8cm}}
\toprule
Ref & Run & Gauge & Length & Sizing basis \\
\midrule
H1 & Pack to protection      & \qty{14}{\awg} & \textcolor{red}{TBD} & Peak bus current \\
H2 & Protection to board     & \qty{14}{\awg} & \textcolor{red}{TBD} & Peak bus current \\
H3 & Board to driver drops (3) & \qty{14}{\awg} & \textcolor{red}{TBD} & Per-axis spin-up current \\
H4 & Motor phase leads (3)   & \textcolor{red}{TBD} & \textcolor{red}{TBD} & Motor phase current \\
H5 & CAN chain segments      & \textcolor{red}{TBD} & stubs $<\qty{10}{\centi\metre}$ & Topology, not current \\
H6 & Encoder SPI runs (3)    & \qtyrange{24}{26}{\awg} & $<\qty{20}{\centi\metre}$ & Length limit governs \\
H7 & Servo drives (3)        & \qtyrange{24}{26}{\awg} & \textcolor{red}{TBD} & Stall current \\
\bottomrule
\end{tabular}
\end{table}

The conductor gauge of the supplied three-conductor CAN cable is not stated by
its supplier, so H5 is an open item: whether stock cables can be used or the bus
harness must be hand-made is decided by that figure, and it affects the
floorplan only marginally but the build time appreciably.

% Figure placeholder: distribution harness drawing -- pack lead, protection
% element, and the three driver drops, with labelling scheme.
% Figure placeholder: as-built harness photograph, labelled and strain-relieved.

\subsection{Board Layout}
\label{app:board_layout}

% Figure placeholder: perfboard floorplan -- bus entry, converter, logic rail,
% bulk capacitance, flight computer, IMU, telemetry module, servo headers and
% the CAN termination network.
% Figure placeholder: board outline drawing with the mounting-hole pattern,
% keep-out heights and connector exit directions delivered to the mechanical
% design as the electronics-mount interface of Table~\ref{tab:mech_interfaces}.

The floorplan is constrained before it is aesthetic: the driver positions follow
from the encoder cable limit, the bus route follows from the driver positions,
and the board is placed around what remains
(Section~\ref{sec:bus_topology}). Rail separation is maintained physically on
the board, not only schematically, so that motor-current return does not share
copper with the sensing ground.

\subsection{Bring-Up and Isolation Checklists}
\label{app:checklists}

These checks are performed on every re-assembly, not only on first build, since
the encoder gaps and connector seating are exactly what mechanical work
disturbs.

\paragraph{Before any power is applied.}
\begin{enumerate}
  \item Isolation: motor bus to chassis, logic rail to motor bus, each rail to
        ground.
  \item Capacitor voltage ratings confirmed by position --- no low-voltage part
        on the pack bus.
  \item Polarity of every fabricated connector verified against
        Table~\ref{tab:connectors}.
  \item Protection element in circuit: E-stop functional, fuse fitted and rated,
        anti-spark path intact.
  \item Encoder air gaps re-measured after mechanical assembly.
\end{enumerate}

\paragraph{First power-on, current-limited bench supply.}
\begin{enumerate}
  \item Current limit verified against a known load before connecting the
        assembly.
  \item Logic rail voltage and ripple under load; converter thermal check.
  \item All drivers enumerate on the bus; encoders read back with no error
        flags; IMU responds.
  \item Bus integrity at the full data rate: sustained soak with all nodes
        addressed, error-frame count logged.
  \item Open-loop wheel motion, wheels contained, one axis at a time.
\end{enumerate}

\paragraph{First battery power-on.} Performed only after the bench sequence
passes in full, with the pack protection verified, cell monitoring active and
the assembly in a fire-safe area. Inrush is measured against the bulk
capacitance on connection.

\subsection{As-Built Measurements}
\label{app:electrical_measured}

Table~\ref{tab:electrical_measured} records the measured electrical figures that
replace the estimates used during design. These are the numbers reported back to
the modelling and control work, and they are the reason the parameter
identification of Section~\ref{sec:paramid} runs on measured rather than assumed
values.

\begin{table}[h]
\centering
\caption{As-built electrical measurements. Each replaces an estimate used
earlier in the design.}
\label{tab:electrical_measured}
\small
\begin{tabular}{L{4.4cm} c c L{4.0cm}}
\toprule
Quantity & Design value & Measured & Replaces / feeds \\
\midrule
Motor torque constant $K_m$    & \qty{0.02513}{\newton\metre\per\ampere} & \textcolor{red}{TBD} & Table~\ref{tab:params}, datasheet value \\
Phase resistance               & \qty{194}{\milli\ohm} & \textcolor{red}{TBD} & Incoming inspection; speed ceiling \\
Wheel inertia $I_w$ (spin-down)& $4.341\times10^{-4}$ & \textcolor{red}{TBD} & Table~\ref{tab:params}; ballast count \\
Viscous and Coulomb friction   & Estimated & \textcolor{red}{TBD} & $C_b$, $C_w$ estimates \\
Command-to-encoder latency     & --- & \textcolor{red}{TBD} & Loop timing, Sec.~\ref{sec:control_sw} \\
Achievable wheel speed on bus  & \qty{6000}{\rpm} & \textcolor{red}{TBD} & $\omega_{\max}$, Sec.~\ref{sec:jumpup_target} \\
Brake torque $\tau_b$          & \qty{5.0}{\newton\metre} (assumed) & \textcolor{red}{TBD} & $\beta$, Sec.~\ref{sec:brake_torque} \\
Logic rail load, worst case    & \textcolor{red}{TBD} & \textcolor{red}{TBD} & Table~\ref{tab:power_budget} \\
Populated board mass           & \textcolor{red}{TBD} & \textcolor{red}{TBD} & Table~\ref{tab:massbudget} \\
Runtime per charge, balancing  & --- & \textcolor{red}{TBD} & Demonstration planning \\
\bottomrule
\end{tabular}
\end{table}

Two rows here gate design decisions rather than merely confirming them. The
brake torque is the unsubstantiated input identified in
Section~\ref{sec:brake_torque}, and it must be measured on the 1-DoF rig before
the wheel geometry is frozen, since $\beta$ propagates straight into the inertia
target. The achievable wheel speed is checked against the design value on a
clamped motor during single-axis bring-up: if the measured ceiling falls below
\qty{6000}{\rpm} on a partly-discharged pack, the conservatism argument of
Section~\ref{sec:jumpup_target} has not been realised and the wheel is
undersized rather than over-margined.
